\par\Needspace{9\baselineskip}\bigskip
\pdfbookmark[2]{First Calculation, [2]}{pallas-III-fragment2}
\begin{center}[2.]\end{center}
\begingroup
\newcommand{\IIILeaderPoint}{.\kern1.1pt}
\newcommand{\IIILeaderUpperLeft}{\rlap{\kern-.05em\raisebox{.5ex}{.}}\IIILeaderPoint}
\newcommand{\IIILeaderUpperRight}{\rlap{\kern.12em\raisebox{.5ex}{.}}\IIILeaderPoint}
\noindent\hspace*{\parindent}[The heliocentric positions of Pallas, that is, the quantities \(r\) and \(v\), were calculated using the system of elements that represents the six oppositions observed so far as closely as possible and that was derived in the Disquisitio de elementis ellipticis Palladis, Art. 14 (Vol. VI, p. 24). Reduced to the mean equinox of 1803 June 30 (the precession in 181 days equals \(24'',\!83\)), this system is as follows.
\begin{center}
\fontsize{10}{12}\selectfont
\setlength{\tabcolsep}{3pt}
\begin{tabular}{@{}p{.68\linewidth}r@{}}
Epoch of the mean longitude for the meridian of Paris \IIILeaderPoint\IIILeaderPoint\IIILeaderPoint\IIILeaderPoint & \(260^{0}\,19'\,30'',\!55\)\\[3pt]
Mean daily tropical motion \IIILeaderPoint\IIILeaderPoint\IIILeaderPoint\IIILeaderPoint\IIILeaderPoint\IIILeaderPoint\IIILeaderPoint\IIILeaderPoint\IIILeaderPoint\IIILeaderPoint\IIILeaderPoint\IIILeaderPoint\IIILeaderPoint\IIILeaderPoint\IIILeaderPoint\IIILeaderPoint\IIILeaderPoint\IIILeaderPoint & \(770'',\!50100\)\\[3pt]
Mean daily sidereal motion \IIILeaderPoint\IIILeaderPoint\IIILeaderPoint\IIILeaderPoint\IIILeaderPoint\IIILeaderPoint\IIILeaderPoint\IIILeaderPoint\IIILeaderPoint\IIILeaderPoint\IIILeaderPoint\IIILeaderPoint\IIILeaderPoint\IIILeaderPoint\IIILeaderPoint\IIILeaderPoint\IIILeaderPoint & \(770'',\!36383\)\\[3pt]
Longitude of perihelion \IIILeaderPoint\IIILeaderPoint\IIILeaderPoint\IIILeaderPoint\IIILeaderPoint\IIILeaderPoint\IIILeaderPoint\IIILeaderPoint\IIILeaderPoint\IIILeaderPoint\IIILeaderPoint\IIILeaderPoint\IIILeaderPoint\IIILeaderPoint\IIILeaderPoint\IIILeaderPoint\IIILeaderPoint\IIILeaderPoint\IIILeaderPoint\IIILeaderPoint\IIILeaderPoint\IIILeaderPoint\IIILeaderPoint\IIILeaderPoint\IIILeaderPoint\IIILeaderPoint\IIILeaderPoint\IIILeaderPoint\IIILeaderPoint & \(121^{0}\,8'\,33'',\!37\)\\[3pt]
Longitude of the ascending node \IIILeaderPoint\IIILeaderPoint\IIILeaderPoint\IIILeaderPoint\IIILeaderPoint\IIILeaderPoint\IIILeaderPoint\IIILeaderPoint\IIILeaderPoint\IIILeaderPoint\IIILeaderPoint\IIILeaderPoint\IIILeaderPoint\IIILeaderPoint\IIILeaderPoint\IIILeaderPoint\IIILeaderPoint\IIILeaderPoint\IIILeaderPoint\IIILeaderPoint\IIILeaderPoint & \(172\quad28\quad37,26\)\\[3pt]
Inclination of the orbit \IIILeaderPoint\IIILeaderUpperLeft\IIILeaderPoint\IIILeaderPoint\IIILeaderPoint\IIILeaderPoint\IIILeaderPoint\IIILeaderPoint\IIILeaderPoint\IIILeaderPoint\IIILeaderPoint\IIILeaderPoint\IIILeaderPoint\IIILeaderPoint\IIILeaderPoint\IIILeaderPoint\IIILeaderPoint\IIILeaderPoint\IIILeaderPoint\IIILeaderPoint\IIILeaderPoint\IIILeaderPoint\IIILeaderPoint\IIILeaderPoint\IIILeaderPoint\IIILeaderPoint\IIILeaderPoint\IIILeaderPoint\IIILeaderPoint\IIILeaderUpperRight\IIILeaderPoint\IIILeaderPoint & \(34\quad37\quad28,35\)\\[3pt]
Eccentricity \(=\sin14^{0}\,9'\,59'',\!79\) \IIILeaderPoint\IIILeaderPoint\IIILeaderPoint\IIILeaderPoint\IIILeaderPoint\IIILeaderPoint\IIILeaderPoint\IIILeaderPoint\IIILeaderPoint\IIILeaderPoint\IIILeaderPoint\IIILeaderPoint\IIILeaderPoint\IIILeaderPoint\IIILeaderPoint\IIILeaderPoint\IIILeaderPoint\IIILeaderPoint\IIILeaderPoint\IIILeaderPoint\IIILeaderPoint\IIILeaderPoint\IIILeaderPoint & \(0,2447424\)\\[3pt]
Logarithm of the semimajor axis \IIILeaderPoint\IIILeaderPoint\IIILeaderPoint\IIILeaderPoint\IIILeaderPoint\IIILeaderPoint\IIILeaderPoint\IIILeaderPoint\IIILeaderPoint\IIILeaderPoint\IIILeaderPoint\IIILeaderPoint\IIILeaderPoint\IIILeaderPoint\IIILeaderPoint\IIILeaderPoint\IIILeaderPoint\IIILeaderPoint\IIILeaderPoint\IIILeaderPoint & \(0,4222071.\)
\end{tabular}
\end{center}
\par
\SourceLeaf{481}
The heliocentric positions of Jupiter needed to calculate the quantities \(r',\,v',\,\beta'\) can be taken from \textsc{Bouvard}’s tables. Since calculating 61 positions from these tables would require very great effort and the utmost precision is unnecessary in the positions of Jupiter, we can use the following device*)\HistoricalUnmarkedNote{[*) See \textsc{Gauss} n \textsc{Olbers} 26 Nov. 1810, p. 415 of this volume.]}, which consists essentially in obtaining a system of elliptic elements for each of the first and second halves of the 3000-day interval under consideration, each representing Jupiter’s motion with sufficient accuracy.
\par
For this we take from \textsc{Bouvard}’s tables the heliocentric positions of Jupiter for the initial epoch 1803 June 30 (Paris) \(=1803+181^{d}\), for the final epoch 1811 September 16 \(=1803+3181^{d}\), and for the epoch at the midpoint, 1807 August 8 \(=1803+1681^{d}\), as given below.
\begin{center}
\setlength{\tabcolsep}{6pt}
\begin{tabular}{r|c|c|r}
 & helioc. longitude & latitude & \(\log r'\)\\[3pt]
\(1803.\ 181^{d}\) & \(188^{0}\,32'\,23'',\!3\) & \(+1^{0}\,18'\,52'',\!6\) & \(0,7367964\)\\[3pt]
\(1681\) & \(307\quad56\quad19,2\) & \(-0\quad38\quad49,7\) & \(0,7061432\)\\[3pt]
\(3181\) & \(82\quad47\quad10,8\) & \(-0\quad21\quad19,8\) & \(0,7081999.\)
\end{tabular}
\end{center}
\par
The reduction to the ecliptic for the three positions is \(+0'',\!2\), \(-23'',\!4\), \(+14'',\!1\), respectively. To obtain the longitudes in the orbit for the three epochs reduced to the mean equinox of 1803 June 30, we take the precession in \(3000^{d}\) to be \(411'',\!5\). The respective longitudes are \(188^{0}\,32'\,23'',\!1\), \(307^{0}\,53'\,16'',\!8\), \(82^{0}\,40'\,5'',\!2\).
\par
From the two positions for \(1803.181^{d}\) and \(1803.1681^{d}\) we can now determine System I of the elements, which serves to calculate Jupiter’s positions for the first half of the interval. The latter two positions yield System II of the elements for the second half of our interval. The heliocentric longitudes and latitudes first give the following by the method set out in Art. 110 of the Theoria motus.
\[
\begin{array}{lll}
\text{System \rlap{\kern-.08em\raisebox{1.7ex}{.}}I.}&\AscendingNode'=98^{0}\,23'\,47''&i'=1\quad18\quad53\\[5pt]
\text{System II.}&\AscendingNode'=98^{0}\,22'\,16''&i'=1\quad18\quad49
\end{array}
\]
\par
The longitudes in the orbit and the radius vectors then give the remaining elements by the method described in the Theoria motus, Arts. 88—96.
\begin{center}
\fontsize{10}{12}\selectfont
\setlength{\tabcolsep}{2pt}
\begin{tabular}{@{}p{.60\linewidth}rr@{}}
 & System I & System II\\[3pt]
\raggedright Epoch of the mean longitude 1807 August 8 \IIILeaderPoint\IIILeaderPoint\IIILeaderPoint\IIILeaderPoint\IIILeaderPoint & \(300^{0}\,45'\,35''\) & \(300^{0}\,38'\,14''\)\\[4pt]
\raggedright Mean daily sidereal motion \IIILeaderPoint\IIILeaderPoint\IIILeaderPoint\IIILeaderPoint\IIILeaderPoint\IIILeaderPoint\IIILeaderPoint\IIILeaderPoint & \(299'',\!1063\) & \(299'',\!1053\)\\[4pt]
\raggedright Longitude of perihelion \IIILeaderPoint\IIILeaderPoint\IIILeaderPoint\IIILeaderPoint\IIILeaderPoint\IIILeaderPoint\IIILeaderPoint\IIILeaderPoint\IIILeaderPoint\IIILeaderPoint\IIILeaderPoint\IIILeaderPoint\IIILeaderPoint\IIILeaderPoint\IIILeaderPoint\IIILeaderPoint\IIILeaderPoint\IIILeaderPoint\IIILeaderPoint\IIILeaderPoint & \(12^{0}\,3'\,14''\) & \(12^{0}\,12'\,17''\)\\[4pt]
\raggedright Logarithm of the eccentricity \IIILeaderPoint\IIILeaderPoint\IIILeaderPoint\IIILeaderPoint\IIILeaderPoint\IIILeaderPoint\IIILeaderPoint\IIILeaderPoint\IIILeaderPoint\IIILeaderPoint\IIILeaderPoint\IIILeaderPoint\IIILeaderPoint\IIILeaderPoint\IIILeaderPoint\IIILeaderPoint & \(8,686027\) & \(8,687916\)\\[4pt]
 & \(e=\sin2^{0}\,46'\,54''\) & \(=\sin2^{0}\,47'\,38''\)\\[4pt]
\raggedright Logarithm of the semimajor axis \IIILeaderPoint\IIILeaderPoint\IIILeaderPoint\IIILeaderPoint\IIILeaderUpperLeft\IIILeaderPoint\IIILeaderPoint\IIILeaderPoint\IIILeaderPoint\IIILeaderPoint\IIILeaderPoint & \(0,716256\) & \(0,716257.\)
\end{tabular}
\end{center}
\par
Evidently, the position of Jupiter for \(1803.1681^{d}\) must be the same when calculated from either system.]
\endgroup
